\section{行列による関係の表現}
\subsection{授業内容}
\subsubsection{科目の中でのこのコマの位置づけ}
線形代数についての基礎的なルールを習得する段階である。今回はより実践的・具体的に，データ行列をどのように線形代数で表現できるかを考える。
データ行列から分散共分散行列，相関行列へと形を変えることを学ぶ。
つづいて線形モデル，とくに従属変数が明確な回帰分析モデルを行列で表現することを見，線形モデルとデザイン行列について考える。
さらに因子分析モデルを行列で表現することを考える。
行列で表現することで，1つの式の中に第一，第二定理の両方を含んだ形で表現できることを理解する。

\subsubsection{コマ主題細目}
\begin{description}
    \item[データの行列表現] 実際に手にするデータセットは，表計算ソフトウェアの画面で見る行列形式の数列であるが，これを記号で表現することで一般的に扱うことができるようになる。添字に気をつけながら要素ごとの表示をすることに加え，行列の計算をこのデータ行列に与えることによって，変数の平均や変数ごとの平均偏差を持った行列が表現できる。平均偏差行列を用いると，行列の積の特徴から分散と共分散を含んだ正方行列が作られることがわかる。また，データを標準化することで，標準化された行列の積が相関行列を表すことになる。このように一般的に表現するために，これまでの行列計算の方法が作られたのだと逆算的に理解すること，加えて行列のサイズに注目しながら，扱うデータの大きさがイメージできるようになることが肝要である。

        \begin{flushright}
            $\to$ \citet{岡太200804}のPp.77-110
        \end{flushright}

    \item[線形モデル] 行列表現の利便性は，データの変換だけにあるのではなく，統計モデルを表現する際にも生きてくる。基礎で学ぶ線形モデルは，基本的にエレメントワイズな表記法であったが，行列を使うことで単回帰も重回帰も同じ式で表現できることがわかる。このように表記の統一性があることが，線形代数の利点である。また統一的な表記にするために，切片項にかかる列を追加するなどの工夫をすることにも注意する。これらの点は，Rなど統計ソフトウェアを扱う上でもヒントになることが多い。

    \item[デザイン行列] 基礎の段階で行った帰無仮説検定は，説明変数が離散変数であったことから，線形モデルの特殊形に過ぎなかったことを再確認する。その上で，先の回帰分析を行列表記にしたように，離散変数で説明する時の係数にかかる行列の形を確認する。この行列はとくにデザイン行列と呼ばれること，また自由度の関係から制約を加えた表現になるが，それがデザイン行列の中でどのように書き表されるかを確認する。

        \begin{flushright}
            $\to$ \citet{Annette_J_Dobson2008-09-08}のPp.41--45にごく簡単な紹介が，\citet{豊田秀樹2000-04-01}のPp.47--62には計画行列として構造方程式の枠組みで説明されている。
        \end{flushright}

        \item[因子分析モデルの行列表現]因子分析モデルはここまでエレメントワイズで表現されていたが，同様に行列表現にするとどのようになるかを確認する。とくに行列のサイズに注目することが重要である。というのも，統計ソフトウェアを使っていると因子得点が表示されないことが少なくないが，行列の形で見ると因子負荷量は項目数$\times$因子数，因子得点は回答者数$\times$因子数になることがより意識されやすいからである。他にも因子分析に関する特徴量が行列のどの要素にはいっているか，また因子分析の定理が行列のどこで表現されているかを確認することが重要である。
        \begin{flushright}
            因子分析の行列的表現については$\to$\citet{芝祐順1979-01}が良書だが，現在は絶版。同様の内容は\citet{kosugi2018}にもある。
        \end{flushright}
\end{description}
\subsubsection{キーワード}
\begin{itemize}
    \item データの行列表現
    \item 分散共分散行列，相関行列
    \item デザイン行列
\end{itemize}
\subsection{授業情報}
\paragraph{コマの展開方法}
講義
\subsubsection{予習・復習課題}
\paragraph{予習}行列の掛け算がメインになってくるので，計算方法並びに計算結果のサイズを確認する方法を見ておこう。
\paragraph{復習}行列表現によって重回帰方程式が1つの形になることを確認する。平均値の差を見るために線形モデルが用いられることを確認する。また因子分析モデルを行列表現すると，一気に2つの定理が1つの式で表現できることを確認する。
